% ===================================================================== Chapter 11
\chapter{Thermo-Structural Analysis}\label{ch:struct}

\section{Structural Model}
The structural model (Table~\ref{tab:smodel}) is a linear static ANSYS Mechanical \cite{ansysmech} model of the Inconel 718 annulus loaded only by the mapped
temperature field. Before any solve, a gate of 61 automated checks compared the written solver inputs with the prepared model section by
section (nodes, elements, temperatures, material, constraints). Each load case solved in about one minute with the sparse direct solver.

%% generated by gen_tables.py - RE-ANALYSIS 2026
\begin{table}[htbp]
\centering
\footnotesize
\caption{Structural model of the baseline (Section 7B)}
\label{tab:smodel}
\begin{tabular}{>{\raggedright\arraybackslash}p{0.30\linewidth}>{\raggedright\arraybackslash}p{0.60\linewidth}}
\toprule
Item & Setting \\
\midrule
Software & ANSYS Mechanical / MAPDL 2026 R1 (Student licence, measured limit 128,000 nodes) \\
Mesh B & SOLID186 quadratic hexahedra, 36 circumferential $\times$ 5 through-wall $\times$ 130 axial (bias 4, 2.13 $\rightarrow$ 8.51 mm); 108,252 nodes / 23,400 elements \\
Element quality & minimum 0.237 (average 0.688); aspect ratio $\le$ 4.88; Jacobian ratio $\le$ 1.205 \\
Thermal load & mapped CFD solid temperature on every node; $\Delta T = T(x,y,z) - T_{ref}$ \\
Reference temperature & 300 K \\
Material & $E(T)$, secant $\alpha(T)$, $\nu$ = 0.294 [ASSUMED]; $S_y(T)$ used only for assessment \\
Analysis & linear static, small deflection, sparse direct solver \\
Load cases & LC1 (free expansion), LC2 (axially restrained, S1), LC2P (LC2 + 443.41 Pa bore pressure) \\
\bottomrule
\end{tabular}
\par\smallskip\parbox{\linewidth}{\footnotesize\raggedright Source: \texttt{MASTER\_PROJECT\_DATA.csv} rows M081--M083; \texttt{PROJECT\_STATE.md} \S13.1 and \S14.1.}
\end{table}


\section{Structural Mesh}
Mesh B uses quadratic 20-node hexahedra (SOLID186), five through the wall (2.0~mm each), 36 around the circumference and 130 along the axis,
graded towards the inlet where the temperature gradients are largest (Figure~\ref{fig:smesh}). With 108,252~nodes it uses 84.6\,\% of the
measured Student limit. Its adequacy is demonstrated by the seven-mesh study of Chapter~\ref{ch:smesh}.

\begin{figure}[htbp]
\centering
\begin{subfigure}[t]{0.58\linewidth}\centering\includegraphics[width=\linewidth,trim=0 80 0 60,clip]{B_mesh_iso.png}\caption{Isometric view}\end{subfigure}\hfill
\begin{subfigure}[t]{0.40\linewidth}\centering\includegraphics[width=\linewidth,trim=250 40 250 40,clip]{B_mesh_end_face.png}\caption{End face: 36 $\times$ 5 elements}\end{subfigure}
\caption[Structural mesh B]{Structural mesh B: 108,252~nodes, 23,400 SOLID186 elements.}
\label{fig:smesh}
\src{\provmech{} The white margins of the export are trimmed for layout.}
\end{figure}

\section{LC1 Free Expansion}
In LC1 the duct is supported in a statically determinate way and can expand freely; the reactions are about $10^{-6}$ N, confirming that no
spurious restraint exists. The maximum total deformation is 1.8443~mm at the outer edge of the outlet face, and the free axial growth
(area-weighted mean over the end face) is 1.8409~mm (Figure~\ref{fig:lc1}). An independent integration of the thermal strain over the same
temperature field gives the same growth to $1.3\times10^{-5}$. The growth is 12.2\,\% below the Section 2 estimate of 2.096~mm because the CFD mean
temperature rise (about 225.5~K) is lower than the analytical one (254.7~K).

\begin{figure}[htbp]
\centering
\begin{subfigure}[t]{0.49\linewidth}\centering\includegraphics[width=\linewidth,height=52mm,keepaspectratio]{LC1_Total_Deformation_iso.png}\caption{Total deformation [m]}\end{subfigure}\hfill
\begin{subfigure}[t]{0.49\linewidth}\centering\includegraphics[width=\linewidth,height=52mm,keepaspectratio]{LC1_Equivalent_Stress_averaged_iso.png}\caption{Von Mises stress [Pa]}\end{subfigure}
\caption[LC1 free expansion]{LC1 free expansion: maximum deformation 1.8443~mm at the outlet end; maximum von Mises stress 24.28~MPa at the bore near the inlet.}
\label{fig:lc1}
\src{\provmech}
\end{figure}

The LC1 stresses come only from the non-uniform temperature. Along most of the duct the bore carries about 18~MPa von Mises (18.20~MPa at
mid-span), the value predicted by the thick-cylinder solution, Equation~\eqref{eq:timoshenko}, evaluated on the CFD field (18.55~MPa). The
maximum, 24.28~MPa at the bore 6.5~mm from the inlet face (427.81~K), is an end effect of the steep axial and radial gradients in the first
slab (Figure~\ref{fig:lcprof}a). This peak lies in the region where the mapped temperatures carry the Fluent node-reconstruction limitation,
so its value is bounded at $\pm$9.8~MPa (Chapter~\ref{ch:mapping}); it drives no conclusion.

\section{LC2 Axially Restrained Expansion}
In LC2 both end faces are held axially. The axial growth is prevented, the maximum total deformation drops to 0.1349~mm (outer surface,
$z \approx$ 244~mm), and the duct carries an almost uniform axial compression (Figure~\ref{fig:lc2}). The end faces react $\pm$548,936.6~N,
balanced to $1.8\times10^{-8}$ N and equal to the analytical restrained force of the same temperature field, Equation~\eqref{eq:restrained},
to $10^{-5}$. The mean axial stress is $-$582.44~MPa.

\begin{figure}[htbp]
\centering
\begin{subfigure}[t]{0.49\linewidth}\centering\includegraphics[width=\linewidth,height=52mm,keepaspectratio]{LC2_Total_Deformation_iso.png}\caption{Total deformation [m]}\end{subfigure}\hfill
\begin{subfigure}[t]{0.49\linewidth}\centering\includegraphics[width=\linewidth,height=52mm,keepaspectratio]{LC2_Equivalent_Stress_averaged_iso.png}\caption{Von Mises stress [Pa]}\end{subfigure}
\caption[LC2 axially restrained expansion (S1 supports)]{LC2 axially restrained expansion (S1 supports): maximum deformation 0.1349~mm; maximum von Mises stress 605.16~MPa at the outer edge of the inlet face.}
\label{fig:lc2}
\src{\provmech}
\end{figure}

\section{Temperature-Dependent Elastic Properties}
Because $E$ decreases and $\alpha$ increases with temperature, the restrained stress $E\alpha\Delta T$ is not proportional to the temperature
rise alone. The analytical model evaluated $E$ and $\alpha$ once at the mean temperature; the FE model evaluates them at every node. Section 2
rescaled to the CFD mean temperature rise gives $-$581.9~MPa, within 0.1\,\% of the FE mean axial stress of $-$582.44~MPa. The difference from
the original analytical $-$657.2~MPa ($-$11.4\,\%) is therefore the temperature basis, not the structural model.

\section{Thermal Stress}
Table~\ref{tab:lcres} summarises both load cases. The LC2 peak von Mises stress of 605.16~MPa occurs at the outer edge of the inlet face
(437.99~K), where the uniform compression combines with the steep temperature gradient of the inlet end (Figure~\ref{fig:lcprof}b). In the interior the von Mises stress is 587--591~MPa, with axial stresses of $-$593.4~MPa at the outer surface and
$-$564.9~MPa at the bore at mid-span. The column load is the end force of 548.94~kN, equivalent to the mean axial
stress of $-$582.44~MPa; the 605.16~MPa is a local peak and not the column load.

%% generated by gen_tables.py - RE-ANALYSIS 2026
\begin{table}[htbp]
\centering
\footnotesize
\caption{Static structural results of the baseline (mesh B)}
\label{tab:lcres}
\begin{tabular}{>{\raggedright\arraybackslash}p{0.42\linewidth}rr}
\toprule
Quantity & LC1 free expansion & LC2 restrained (S1) \\
\midrule
Maximum total deformation [mm] & 1.8443 & 0.1349 \\
Free axial growth $\Delta L$ [mm] & 1.8409 & -- \\
Maximum von Mises stress [MPa] & 24.282 & 605.161 \\
Location of the maximum & bore, $z$ = 6.5 mm & outer edge of the inlet face \\
Temperature at the maximum [K] & 427.81 & 437.99 \\
Mid-span bore von Mises [MPa] & 18.20 & -- \\
Axial end reaction [N] & $\approx 0$ ($9.0\times10^{-7}$) & 548,936.6 \\
Mean axial stress [MPa] & $\approx 0$ & $-$582.440 \\
Pressure effect LC2P $-$ LC2 on the peak [Pa] & -- & +45 \\
\bottomrule
\end{tabular}
\par\smallskip\parbox{\linewidth}{\footnotesize\raggedright Source: \texttt{MASTER\_PROJECT\_DATA.csv} rows M084--M097 (VERIFIED from the MAPDL nodal and reaction tables).}
\end{table}



The ratio of restrained to gradient stress is 32 at mid-span and 25 between the peaks: the axial restraint, not the through-wall temperature
gradient, dominates the thermal stress. This confirms in kind the Section 2 prediction of a factor of about 40.

\begin{figure}[htbp]
\centering
\begin{subfigure}[t]{0.45\linewidth}\centering\includegraphics[width=\linewidth,height=78mm,keepaspectratio]{F7B_01_LC1_axial_profiles.png}\caption{LC1: temperature, axial displacement and von Mises stress}\end{subfigure}\hfill
\begin{subfigure}[t]{0.53\linewidth}\centering\includegraphics[width=\linewidth,height=78mm,keepaspectratio]{F7B_02_LC2_axial_profiles.png}\caption{LC2: temperature, axial displacement, von Mises and axial stress}\end{subfigure}
\caption{Axial profiles (circumferential means at the bore and outer surface) along the duct for LC1 and LC2.}
\label{fig:lcprof}
\src{\provdata}
\end{figure}


\section{Deformation}
The LC1 growth is distributed along the duct in proportion to the local temperature rise, so the axial displacement grows faster towards the
hot outlet (Figure~\ref{fig:lcprof}a). In LC2 the axial displacement is small and mainly an internal redistribution: the axial displacement reaches
at most 0.109~mm, and the maximum radial displacement is 0.0897~mm. The small LC2 deformation hides a large stored strain energy, which is
the reason the restrained duct must also be assessed for stability.

\section{Reaction Forces}
In LC2 the axial force is constant along the length and equal to the end
reaction; in LC1 the section force computed from the corner-node stresses is about 125~N, i.e.\ zero to within about 0.02\,\% of the LC2 force, as it must be for a free body. The three
mid-span hoop nodes of LC2 react about $10^{-7}$ N; they only suppress rigid-body rotation.


\section{Pressure-Load Effect}
The fluid pressure acts on the bore. Load case LC2P applied the highest CFD wall pressure, 443.41~Pa gauge, uniformly over the whole bore as
a bound (atmospheric pressure acts on all surfaces and was not applied). The peak von Mises stress changed from 605.160873 to 605.160918~MPa,
i.e.\ by +45~Pa or $7.4\times10^{-6}$\,\%; a Lam\'e check of the pressure-induced stress agreed within 0.3\,\%.
Pressure loading is negligible compared with the thermal loading and was not re-applied in the parametric cases, whose pressure drops are
of the same order.


The first-yield assessment of the LC2 static state is given in Table~\ref{tab:yield}. The local yield strength at the critical node
(164.8\,\textdegree C) is 1047.0~MPa, giving a utilisation of 0.578 and a first-yield load factor of 1.730: the static state is elastic, and the
thermal load would have to rise by 73\,\% before the first point yields. This is a statement about the static stress state only. It is not a
structural margin, because the restrained duct becomes unstable before that load level is reached under the S1 supports
(Chapter~\ref{ch:buckling}).

%% generated by gen_tables.py - RE-ANALYSIS 2026
\begin{table}[htbp]
\centering
\footnotesize
\caption{First-yield assessment of the LC2 static state (not a structural margin; see Chapter~\ref{ch:buckling})}
\label{tab:yield}
\begin{tabular}{>{\raggedright\arraybackslash}p{0.55\linewidth}r}
\toprule
Quantity & Value \\
\midrule
Critical node temperature [K] & 437.99 \\
Local yield strength $S_y(T)$ at the critical node [MPa] & 1047.0 \\
Yield utilisation $\sigma_{vM}/S_y(T)$ & 0.5780 \\
First-yield load factor (1/utilisation) & 1.730 \\
First-yield margin $S_y/\sigma_{vM} - 1$ & 0.730 \\
LC1 utilisation & 0.023 \\
\bottomrule
\end{tabular}
\par\smallskip\parbox{\linewidth}{\footnotesize\raggedright Source: \texttt{MASTER\_PROJECT\_DATA.csv} rows M093--M096; \texttt{STRUCTURAL\_RESULTS.md}. $S_y(T)$ interpolated linearly in the VDM table at the node temperature; the Engineering Data scalar of 1020~MPa is a lower bound and is not used.}
\end{table}

